% Unit 045 English translation of chapter3.tex lines 2714--2935.
% Source: Methods of Algebra, Volume 2, commit
% 9a5803ff77dd3257484cb177f851a73770a59dd3 (CC BY 4.0).
% stable-unit-id: o014.aljabr2.chapter3.ext-tor
% Coverage: Section 3.14 in full.

% segment-id: o014.li2.chapter3.ext-tor.h001
\section{Example: \texorpdfstring{$\Ext$}{Ext} and \texorpdfstring{$\Tor$}{Tor}}\label{sec:Ext-Tor}
\hypertarget{o014.aljabr2.chapter3.ext-tor}{ }

% segment-id: o014.li2.chapter3.ext-tor.p001
We shall study the derived functors of $\Hom$ and $\otimes$. It is useful to
introduce a few auxiliary concepts first.

% segment-id: o014.li2.chapter3.ext-tor.de001
\begin{definition}\label{def:balanced-functor}
\index{bifunctor!balanced}
Let $\mathcal{A}_1$, $\mathcal{A}_2$, and $\mathcal{B}$ be abelian categories,
and let the bifunctor $F:\mathcal{A}_1\times\mathcal{A}_2\to\mathcal{B}$
(Convention~\ref{con:bifunctor}) be left exact in both variables. If, for
every injective object $I_i\in\Obj(\mathcal{A}_i)$ ($i=1,2$), both functors
\[
F(I_1,\cdot):\mathcal{A}_2\to\mathcal{B}\quad\text{and}\quad
F(\cdot,I_2):\mathcal{A}_1\to\mathcal{B}
\]
are exact, then $F$ is called \emph{balanced}.

Dually, if $F$ is right exact in both variables, and if for every projective
object $P_i\in\Obj(\mathcal{A}_i)$ the functors $F(P_1,\cdot)$ and
$F(\cdot,P_2)$ are both exact, then $F$ is again called balanced.
\end{definition}

% segment-id: o014.li2.chapter3.ext-tor.p002
By duality, the following theorem is stated only for bifunctors that are left
exact in both variables.

% segment-id: o014.li2.chapter3.ext-tor.th001
\begin{theorem}\label{prop:balanced-primer}
\index[sym1]{RInF@$\mathrm{R}_{\mathrm{I}}^n F$, $\mathrm{R}_{\mathrm{II}}^n F$}
Let $\mathcal{A}_1$, $\mathcal{A}_2$, and $\mathcal{B}$ be abelian
categories, and suppose that $\mathcal{A}_1$ and $\mathcal{A}_2$ have enough
injective objects. Let the bifunctor
$F:\mathcal{A}_1\times\mathcal{A}_2\to\mathcal{B}$ be left exact in both
variables. We may therefore take right derived functors in either variable.
For fixed $X_i\in\Obj(\mathcal{A}_i)$, $i=1,2$, write
\[
\mathrm{R}_{\mathrm{I}}^n F(\cdot,X_2):\mathcal{A}_1\to\mathcal{B},\qquad
\mathrm{R}_{\mathrm{II}}^n F(X_1,\cdot):\mathcal{A}_2\to\mathcal{B}.
\]
If $F$ is balanced, then there is a canonical isomorphism of bifunctors
\[
\mathrm{R}_{\mathrm{I}}^nF(X_1,X_2)\simeq
\mathrm{R}_{\mathrm{II}}^nF(X_1,X_2).
\]
\end{theorem}

% segment-id: o014.li2.chapter3.ext-tor.pf001
\begin{proof}
The following argument relies on the theory of double complexes; see
\S\ref{sec:double-cplx} and \S\ref{sec:double-cplx-coh}. For $i=1,2$, choose
an injective resolution
$0\to X_i\to I_i^0\to I_i^1\to\cdots$; define $I_i^n:=0$ for $n<0$.
This defines the double complex $F(I_1^\bullet,I_2^\bullet)$.

On the other hand, form the complexes $F(X_1,I_2^\bullet)$ and
$F(I_1^\bullet,X_2)$. Regard them as double complexes concentrated,
respectively, on the vertical axis $(0,\bullet)$ and the horizontal axis
$(\bullet,0)$. Since
$X_i\rightiso\Ker[I_i^0\to I_i^1]\hookrightarrow I_i^0$, there are morphisms
in $\cate{C}^2_f(\mathcal{B})$
\[
F(X_1,I_2^\bullet)\xrightarrow{\lambda}
F(I_1^\bullet,I_2^\bullet)\xleftarrow{\rho}
F(I_1^\bullet,X_2).
\]
We shall show that $\tot(\lambda)$ and $\tot(\rho)$ are quasi-isomorphisms.

First consider $\lambda$. Take horizontal cohomology $\Hm_{\mathrm{I}}$ on
both sides; this leaves $F(X_1,I_2^\bullet)$, which lies on the vertical
axis, unchanged. Since $F(\cdot,I_2^q)$ is exact for every $q$, we obtain
\[
\left(\Hm_{\mathrm{I}}F(I_1^\bullet,I_2^\bullet)\right)^{p,q}
=\begin{cases}
0,&p\neq0,\\
\Ker\left[F(I_1^0,I_2^q)\to F(I_1^1,I_2^q)\right]
=F(X_1,I_2^q),&p=0.
\end{cases}
\]
Thus $\Hm_{\mathrm{I}}(\lambda)$ is an isomorphism from
$\Hm_{\mathrm{I}}F(X_1,I_2^\bullet)$ to
$\Hm_{\mathrm{I}}F(I_1^\bullet,I_2^\bullet)$, so
$\Hm_{\mathrm{II}}\Hm_{\mathrm{I}}(\lambda)$ is also an isomorphism.
Theorem~\ref{prop:double-cplx-tot} now shows that $\tot(\lambda)$ is a
quasi-isomorphism.

For $\rho$, the same argument shows that
$\Hm_{\mathrm{I}}\Hm_{\mathrm{II}}(\rho)$ is an isomorphism;
Theorem~\ref{prop:double-cplx-tot} then implies that $\tot(\rho)$ is also a
quasi-isomorphism.

The total complex of $F(X_1,I_2^\bullet)$ is still that complex itself, and
taking $\Hm^n$ gives $\mathrm{R}_{\mathrm{II}}^nF(X_1,X_2)$. Likewise,
taking $\Hm^n$ of $F(I_1^\bullet,X_2)$, or of its total complex, gives
$\mathrm{R}_{\mathrm{I}}^nF(X_1,X_2)$. Using
$\tot(F(I_1^\bullet,I_2^\bullet))$ as an intermediary yields
\[
\mathrm{R}_{\mathrm{I}}^nF(X_1,X_2)\simeq
\mathrm{R}_{\mathrm{II}}^nF(X_1,X_2).
\]
Because every morphism $X_i\to Y_i$ lifts to a morphism between injective
resolutions, uniquely up to homotopy, what we have obtained is in fact an
isomorphism of bifunctors
$\mathrm{R}_{\mathrm{I}}^nF\simeq\mathrm{R}_{\mathrm{II}}^nF$.
\end{proof}

% segment-id: o014.li2.chapter3.ext-tor.rem001
\begin{remark}\label{rem:balanced-primer}
The argument actually gives a method for resolving both variables
simultaneously and then deriving $F$ in a balanced way by means of the total
complex of the double complex $F(I_1^\bullet,I_2^\bullet)$.
\end{remark}

% segment-id: o014.li2.chapter3.ext-tor.p003
Now let $\mathcal{A}$ be an abelian category and consider the bifunctor
$\Hom=\Hom_{\mathcal{A}}:\mathcal{A}^{\opp}\times\mathcal{A}\to\cate{Ab}$.
Proposition~\ref{prop:Hom-left-exact} says that $\Hom$ is left exact in both
variables. For every $n\in\Z$, make the following definitions.

% segment-id: o014.li2.chapter3.ext-tor.l001
\begin{itemize}
% segment-id: o014.li2.chapter3.ext-tor.i001
\item If $\mathcal{A}$ has enough projective objects, define
$\Ext^n_{\mathcal{A},\mathrm{I}}(X,Y):=\mathrm{R}^n_{\mathrm{I}}\Hom(X,Y)$,
computed using a projective resolution of $X$ in $\mathcal{A}$.
% segment-id: o014.li2.chapter3.ext-tor.i002
\item If $\mathcal{A}$ has enough injective objects, define
$\Ext^n_{\mathcal{A},\mathrm{II}}(X,Y):=\mathrm{R}^n_{\mathrm{II}}\Hom(X,Y)$,
computed using an injective resolution of $Y$ in $\mathcal{A}$.
\end{itemize}

% segment-id: o014.li2.chapter3.ext-tor.dp001
\begin{definition-proposition}[The $\Ext$ functor]\label{def:Ext-classical}
\index{Ext functor@$\Ext$ functor}
\index[sym1]{Extn@$\Ext^n$}
The bifunctor $\Hom_{\mathcal{A}}$ is balanced. Consequently, whenever the
abelian category $\mathcal{A}$ has enough injective and projective objects,
there is, for every $n\in\Z$, a bifunctor
\[
\Ext^n=\Ext^n_{\mathcal{A}}:\mathcal{A}^{\opp}\times\mathcal{A}\to\cate{Ab}
\]
such that
$\Ext^n_{\mathcal{A}}(X,Y):=\Ext^n_{\mathrm{I}}(X,Y)\simeq
\Ext^n_{\mathrm{II}}(X,Y)$ (canonically).
\end{definition-proposition}

% segment-id: o014.li2.chapter3.ext-tor.pf002
\begin{proof}
Balancedness is an immediate consequence of the definitions of injective and
projective objects, so Theorem~\ref{prop:balanced-primer} applies.
\end{proof}

% segment-id: o014.li2.chapter3.ext-tor.p004
If $R$ is a ring and $\mathcal{A}=R\dcate{Mod}$ or $\cated{Mod}R$, we write
$\Ext^n_{\mathcal{A}}$ as $\Ext^n_R$.

% segment-id: o014.li2.chapter3.ext-tor.p005
Observe that $\Ext^0=\Hom$, while $(\Ext^n)_{n\geq0}$ has, in the variables
$X$ and $Y$ respectively, long exact sequences of the form
% segment-id: o014.li2.chapter3.ext-tor.q001
\begin{gather*}
\cdots\to\Ext^{n-1}(X',Y)\xrightarrow{\delta^{n-1}}\Ext^n(X'',Y)\to
\Ext^n(X,Y)\to\Ext^n(X',Y)\to\cdots,\\
\cdots\to\Ext^{n-1}(X,Y'')\xrightarrow{\delta^{n-1}}\Ext^n(X,Y')\to
\Ext^n(X,Y)\to\Ext^n(X,Y'')\to\cdots,
\end{gather*}
where $0\to X'\to X\to X''\to0$ and
$0\to Y'\to Y\to Y''\to0$ are given short exact sequences.

% segment-id: o014.li2.chapter3.ext-tor.p006
If $\mathcal{A}$ is a $\Bbbk$-linear abelian category, the values of $\Ext^n$
may naturally be regarded as objects of $\Bbbk\dcate{Mod}$.

% segment-id: o014.li2.chapter3.ext-tor.rem002
\begin{remark}
The symbol $\Ext$ stands for “extension,” because the elements of
$\Ext^1(X,Y)$ correspond bijectively to the isomorphism classes of extensions
$0\to Y\to E\to X\to0$ (that is, short exact sequences); the higher
$\Ext^n$ have a similar interpretation. This fact is best understood through
the derived category; see \S\sourcecrossref{sec:Hom-Ext}{4.5}.
\end{remark}

% segment-id: o014.li2.chapter3.ext-tor.p007
The characterization of exact functors in
Corollary~\ref{prop:obstruction-exactness} immediately gives the following.

% segment-id: o014.li2.chapter3.ext-tor.l002
\begin{itemize}
% segment-id: o014.li2.chapter3.ext-tor.i003
\item An object $I$ is injective $\iff \Ext^1(\cdot,I)=0\iff
\Ext^{\geq1}(\cdot,I)=0$.
% segment-id: o014.li2.chapter3.ext-tor.i004
\item An object $P$ is projective $\iff \Ext^1(P,\cdot)=0\iff
\Ext^{\geq1}(P,\cdot)=0$.
\end{itemize}

% segment-id: o014.li2.chapter3.ext-tor.ex001
\begin{example}[$\HHm^n$ as $\Ext^n$]\label{eg:HH-Ext}
\index[sym1]{HH}
Return to the basic setting for Hochschild homology and cohomology from
\S\ref{sec:HH}: $\Bbbk$ is a commutative ring and $R$ is a
$\Bbbk$-algebra; now assume in addition that $R$ is projective as a
$\Bbbk$-module. Write $\otimes:=\otimes_{\Bbbk}$. Since projective modules
are direct summands of free modules, $R^{\otimes n}$ is still projective as a
$\Bbbk$-module.
\footnote{Alternatively, use the tensor adjunction
\cite[Theorem 6.6.5]{Li1}.}
Regard $R\otimes R^{\otimes n}\otimes R$ as a left $R^e$-module, where
$R^e:=R\otimes R^{\opp}$. Since
$R\otimes(\cdot)\otimes R:\Bbbk\dcate{Mod}\to R^e\dcate{Mod}$ is left
adjoint to the forgetful functor \cite[Corollary 6.6.8]{Li1}, it carries
projective objects to projective objects
(Proposition~\ref{prop:adjoint-injective-projective}). Thus
Lemma~\ref{prop:HH-bar-exactness} shows that the bar complex $\mathsf{B}R$
gives a projective resolution of $R$ as a left $R^e$-module,
\begin{equation}\label{eqn:bar-proj-resolution}
\cdots\to\mathsf{B}_1R\to\mathsf{B}_0R\to R\to0.
\end{equation}
Now let $M$ be an $(R,R)$-bimodule. Substituting the original definition of
Hochschild cohomology from Definition~\ref{def:HH} gives
\[
\HHm^n(M)=\Hm^n\left(\Hom_{R^e}(\mathsf{B}R,M)\right)
=\Ext^n_{R^e}(R,M).
\]
\end{example}

% segment-id: o014.li2.chapter3.ext-tor.p008
We next turn to the $\Tor$ functor. Let $R$ be a ring. The module categories
$R\dcate{Mod}$ (and $\cated{Mod}R$) have enough projective
objects\footnote{Indeed, free modules are always projective; see
\cite[Example 6.9.6]{Li1}.}. The tensor product gives a bifunctor
% segment-id: o014.li2.chapter3.ext-tor.q002
\[
\otimes_R:\cated{Mod}R\times R\dcate{Mod}\to\cate{Ab},
\]
which is right exact in either variable (see
\cite[Proposition 6.9.2]{Li1}). Apply the right-exact version of the preceding
theory to define the left derived functors in both variables,
$\mathrm{L}_{\mathrm{I},n}\otimes_R$ and
$\mathrm{L}_{\mathrm{II},n}\otimes_R$; these can be nonzero only for
$n\geq0$. Recall also the definition of a flat module
\cite[Definition 6.9.4]{Li1}.

% segment-id: o014.li2.chapter3.ext-tor.dp002
\begin{definition-proposition}[The $\Tor$ functor]\label{def:Tor-classical}
\index{Tor functor@$\Tor$ functor}
\index[sym1]{Torn@$\Tor^R_n$}
For every ring $R$, the bifunctor $\otimes_R$ is balanced. Hence, for every
$n\in\Z$, one may define a bifunctor
\[
\Tor_n=\Tor_n^R:\cated{Mod}R\times R\dcate{Mod}\to\cate{Ab}
\]
by
$\Tor_n^R(X,Y):=(\mathrm{L}_{\mathrm{I},n}\otimes_R)(X,Y)\simeq
(\mathrm{L}_{\mathrm{II},n}\otimes_R)(X,Y)$ (canonically).
\end{definition-proposition}

% segment-id: o014.li2.chapter3.ext-tor.pf003
\begin{proof}
Balancedness follows from the fact that every projective module is flat
\cite[Corollary 6.9.9]{Li1}. Apply the left-derived-functor version of
Theorem~\ref{prop:balanced-primer}.
\end{proof}

% segment-id: o014.li2.chapter3.ext-tor.p009
We have $\Tor^R_0(X,Y)=X\dotimes{R}Y$, while $(\Tor_n)_{n\geq0}$ has the
following long exact sequences:
% segment-id: o014.li2.chapter3.ext-tor.q003
\begin{gather*}
\cdots\to\Tor^R_{n+1}(X'',Y)\xrightarrow{\partial_{n+1}}\Tor^R_n(X',Y)\to
\Tor^R_n(X,Y)\to\Tor^R_n(X'',Y)\to\cdots,\\
\cdots\to\Tor^R_{n+1}(X,Y'')\xrightarrow{\partial_{n+1}}\Tor^R_n(X,Y')\to
\Tor^R_n(X,Y)\to\Tor^R_n(X,Y'')\to\cdots,
\end{gather*}
where $0\to X'\to X\to X''\to0$ and
$0\to Y'\to Y\to Y''\to0$ are given short exact sequences.

% segment-id: o014.li2.chapter3.ext-tor.p010
As in the case of $\Ext$, we obtain the following characterizations.

% segment-id: o014.li2.chapter3.ext-tor.l003
\begin{itemize}
% segment-id: o014.li2.chapter3.ext-tor.i005
\item A right $R$-module $M$ is flat $\iff \Tor^R_1(M,\cdot)=0\iff
\Tor^R_{\geq1}(M,\cdot)=0$.
% segment-id: o014.li2.chapter3.ext-tor.i006
\item A left $R$-module $N$ is flat $\iff \Tor^R_1(\cdot,N)=0\iff
\Tor^R_{\geq1}(\cdot,N)=0$.
\end{itemize}

% segment-id: o014.li2.chapter3.ext-tor.p011
Let $M$ be an $R$-module. In analogy with the definition of a projective
resolution, a \emph{flat resolution}\index{flat resolution} of $M$ is an
exact sequence of the form
% segment-id: o014.li2.chapter3.ext-tor.q004
\[
\cdots\to F^2\to F^1\to F^0\to M\to0,\qquad
\text{each $F^n$ is a flat $R$-module}.
\]

% segment-id: o014.li2.chapter3.ext-tor.p012
By Corollary~\ref{prop:acyclic-resolution} and the preceding observation,
$\Tor^R_n(X,Y)$ can also be computed using a flat resolution of $X$ or $Y$.
This is the practical technique for computing $\Tor$, whereas projective
resolutions play a chiefly theoretical role.

% segment-id: o014.li2.chapter3.ext-tor.p013
Such computations can naturally be balanced as well. Choose arbitrary flat
resolutions $\cdots\to P_1\to P_0\to X\to0$ and
$\cdots\to Q_1\to Q_0\to Y\to0$. Form the double complex, in the sense of
chain complexes, $P_\bullet\dotimes{R}Q_\bullet$, and denote its total
complex by $P\otimes Q$. The homological version of
Remark~\ref{rem:balanced-primer} gives
$\Tor_n^R(X,Y)\simeq\Hm_n(P\otimes Q)$; the only properties needed are the
exactness of $P_p\dotimes{R}(\cdot)$ and
$(\cdot)\dotimes{R}Q_q$, precisely the properties guaranteed by flatness.

% segment-id: o014.li2.chapter3.ext-tor.rem003
\begin{remark}[Bimodule structure]\label{rem:ExtTor-bimodule}
Let $A$ and $B$ be rings, let $X$ be an $(A,R)$-bimodule, and let $Y$ be an
$(R,B)$-bimodule. Since $\Tor_n^R$ is a bifunctor, $\Tor_n^R(X,Y)$ carries a
canonical $(A,B)$-bimodule structure. For example, consider left
multiplication by $A$: every $a\in A$ gives an endomorphism of $X$ as a
right $R$-module and hence induces an endomorphism $L_a$ of
$\Tor_n^R(X,Y)$; one has $L_aL_{a'}=L_{aa'}$ and $L_1=\identity$. Right
multiplication by $B$ is handled similarly in the second variable. For
$n=0$, this recovers the natural bimodule structure on $X\dotimes{R}Y$.

If $Y$, as a left $R$-module, is supplied with a flat resolution
$\cdots\to P_0\to Y\to0$, the left $A$-action can be read directly from the
complex $X\dotimes{R}P_\bullet$: it comes from left multiplication by $A$ on
$X$. Right multiplication is treated in the same way.

As a special case, if a ring homomorphism $R\to S$ is fixed, then
$\Tor_n^R(\cdot,S)$ (or $\Tor_n^R(S,\cdot)$) naturally becomes a functor
$\cated{Mod}R\to\cated{Mod}S$ (or
$R\dcate{Mod}\to S\dcate{Mod}$).

A similar argument shows that, when $X$ is an $(R,A)$-bimodule and $Y$ is an
$(R,B)$-bimodule, every $\Ext_R^n(X,Y)$ naturally carries an
$(A,B)$-bimodule structure.
\end{remark}

% segment-id: o014.li2.chapter3.ext-tor.p014
For a commutative ring $R$, there is no need to distinguish left from right
$R$-modules; the corresponding category is denoted by $R\dcate{Mod}$.
Every $R$-module is then naturally an $(R,R)$-bimodule, so $\Tor_n^R$ becomes
a bifunctor $(R\dcate{Mod})^2\to R\dcate{Mod}$.

% segment-id: o014.li2.chapter3.ext-tor.pr001
\begin{proposition}\label{prop:Tor-symmetry}
For a commutative ring $R$, there is a canonical isomorphism of bifunctors
\[
\Tor_n^R(X,Y)\simeq\Tor_n^R(Y,X)
\]
for every $n\in\Z$.
\end{proposition}

% segment-id: o014.li2.chapter3.ext-tor.pf004
\begin{proof}
Choose flat resolutions $P_\bullet$ and $Q_\bullet$ of $X$ and $Y$,
respectively. Write $P\otimes Q$ for the total complex of the double complex
$P_\bullet\dotimes{R}Q_\bullet$, and define $Q\otimes P$ similarly. The
symmetry constraint for the tensor product
\cite[Proposition 6.5.14]{Li1} gives canonical isomorphisms
% segment-id: o014.li2.chapter3.ext-tor.q005
\[
P_p\dotimes{R}Q_q\simeq Q_q\dotimes{R}P_p,\qquad p,q\in\Z.
\]
In other words,
$P_\bullet\dotimes{R}Q_\bullet\simeq
\mathrm{swap}(Q_\bullet\dotimes{R}P_\bullet)$. Apply
Proposition~\ref{prop:double-cplx-swap} to obtain a canonical isomorphism of
total complexes $P\otimes Q\simeq Q\otimes P$. This proves the assertion.
\end{proof}

% segment-id: o014.li2.chapter3.ext-tor.rem004
\begin{remark}
The symbol $\Tor$ stands for “torsion,” as suggested by the following
example. Let $\mathfrak{a}$ be a left ideal of a ring $R$, and let $X$ be any
right $R$-module. Since the free left module $R$ is flat, apply
Proposition~\ref{prop:dimension-shifting} to
$0\to\mathfrak{a}\to R\to R/\mathfrak{a}\to0$ to shift dimensions. This
gives
% segment-id: o014.li2.chapter3.ext-tor.q006
\begin{align*}
\Tor^R_1(X,R/\mathfrak{a}) &\simeq \Ker\left[
X\dotimes{R}\mathfrak{a}\to X\dotimes{R}R\simeq X\right]\\
&\simeq \left\{\sum_i x_i\otimes a_i\in X\dotimes{R}\mathfrak{a}:
\sum_i x_i a_i=0\right\},\\
\Tor^R_n(X,R/\mathfrak{a})&\simeq\Tor^R_{n-1}(X,\mathfrak{a}),\qquad(n>1).
\end{align*}
If, moreover, $t\in R$ is not a right zero divisor and we take
$\mathfrak{a}:=Rt\leftiso R$ (mapping $r$ to $rt$), then
$\Tor^R_1(X,R/Rt)\simeq\{x\in X:xt=0\}$, the $t$-torsion submodule of $X$;
for $n>1$, $\Tor^R_n(X,R/Rt)=0$.
\end{remark}

% segment-id: o014.li2.chapter3.ext-tor.ex002
\begin{example}[$\HHm_n$ as $\Tor_n$]\label{eg:HH-Tor}
\index[sym1]{HH}
Continue from the opening setup of Example~\ref{eg:HH-Ext}, but weaken the
condition on $R$ by requiring only that $R$ be flat as a $\Bbbk$-module.
Then $R^{\otimes n}$ is flat as well. The associativity constraint
$(\cdot)\dotimes{R^e}(R\otimes R^{\otimes n}\otimes R)\simeq
(\cdot)\otimes R^{\otimes n}$ shows that
$R\otimes R^{\otimes n}\otimes R$ is a flat left $R^e$-module. Thus
$\mathsf{B}R$ is a flat resolution of $R$ as a left $R^e$-module, as in
\eqref{eqn:bar-proj-resolution}.

For every $(R,R)$-bimodule $M$, the original definition of Hochschild
homology from Definition~\ref{def:HH} gives
\[
\HHm_n(M)=\Hm_n\left(M\dotimes{R^e}\mathsf{B}R\right)
=\Tor_n^{R^e}(M,R).
\]
\end{example}

% segment-id: o014.li2.chapter3.ext-tor.p015
We can now combine Examples~\ref{eg:HH-Ext} and \ref{eg:HH-Tor} to compute
further examples of Hochschild homology and cohomology. First take the
polynomial algebra $R=\Bbbk[t]$ and identify
$R^e=R\otimes R^{\opp}$ with $\Bbbk[x,y]$. The scalar action on $R$ as an
$R^e$-module is $h(x,y)\cdot f(t)=h(t,t)f(t)$, and $R$ has the free
resolution
% segment-id: o014.li2.chapter3.ext-tor.q007
\[
0\to R^e\xrightarrow{\text{multiplication by }x-y}R^e
\xrightarrow{f\mapsto f(t,t)}R\to0.
\]
Computing $\Ext^n_{R^e}(R,\cdot)$ and $\Tor_n^{R^e}(\cdot,R)$ from this
resolution immediately gives $\HHm_n(R)\simeq R\simeq\HHm^n(R)$ for
$n\in\{0,1\}$, while $\HHm_n(R)=0=\HHm^n(R)$ for $n>1$.

% segment-id: o014.li2.chapter3.ext-tor.p016
Next take $R=\Bbbk[t]/(t^2)$, so that
$R^e=\Bbbk[x,y]/(x^2,y^2)$. As an $R^e$-module, $R$ has a free resolution
of period $2$,
% segment-id: o014.li2.chapter3.ext-tor.q008
\[
\cdots\xrightarrow{x+y}R^e\xrightarrow{x-y}R^e\xrightarrow{x+y}R^e
\xrightarrow{\text{multiplication by }x-y}R^e\xrightarrow{f\mapsto f(t,t)}R\to0,
\]
whose exactness is left to the reader to verify directly. For every
$R$-module $M$, applying $M\dotimes{R^e}(\cdot)$ and
$\Hom_{R^e}(\cdot,M)$ gives, respectively,
% segment-id: o014.li2.chapter3.ext-tor.q009
\begin{gather*}
\cdots\xrightarrow{2t}M\xrightarrow{0}M\xrightarrow{2t}M\xrightarrow{0}M
\to M\dotimes{R^e}R\to0,\\
0\to\Hom_{R^e}(R,M)\to M\xrightarrow{0}M\xrightarrow{2t}M
\xrightarrow{0}M\xrightarrow{2t}\cdots.
\end{gather*}
Thus $\HHm_k(M)$ and $\HHm^k(M)$ depend only on $k\bmod2$; this periodicity
might be difficult to see from the definitions alone.

If $\Bbbk$ is a field, $\HHm_n(M)$ and $\HHm^n(M)$ can be read off from
these complexes, but the answer depends on $\mathrm{char}(\Bbbk)$. The
exercises will treat the general case $R=\Bbbk[t]/(t^n)$.

% segment-id: o014.li2.chapter3.ext-tor.p017
The following result is often used in algebraic topology and also provides a
small check on the preceding material; it concerns chain complexes and their
homology (Remark~\ref{rem:cochain-vs-chain}). Let
$C=(C_\bullet,d^C_\bullet)$ (respectively,
$D=(D_\bullet,d^D_\bullet)$) be a chain complex of right $R$-modules
(respectively, left $R$-modules). Form the double chain complex
$C_\bullet\otimes D_\bullet$ and denote its total complex by
% segment-id: o014.li2.chapter3.ext-tor.q010
\[
C\otimes D:=\tot_{\oplus}\left(C_\bullet\otimes D_\bullet\right);
\]
the subscript $R$ on the tensor product is omitted. Concretely, the
degree-$n$ term of the total complex $C\otimes D$ is
$\bigoplus_{p+q=n}C_p\otimes D_q$, while the restriction of $d_n$ to
$C_p\otimes D_q$ is
$d^C_p\otimes\identity+(-1)^p\identity\otimes d^D_q$. This gives a canonical
homomorphism
\[
\kappa:\bigoplus_{p+q=n}\Hm_p(C)\otimes\Hm_q(D)\to
\Hm_n(C\otimes D).
\]

% segment-id: o014.li2.chapter3.ext-tor.th002
\begin{theorem}[Künneth theorem for homology]\label{prop:Kunneth-homology}
\index{Kunneth@Künneth theorem}
Let $C,D$ be chain complexes as above. If every $C_p$ and every
$\Image(d^C_p)$ is a flat right $R$-module, then there is a canonical short
exact sequence
% segment-id: o014.li2.chapter3.ext-tor.g001
\[
\begin{tikzcd}[column sep=small, every cell/.append style={font = \small}]
0 \arrow[r] & \displaystyle\bigoplus_{p+q=n} \Hm_p(C) \otimes \Hm_q(D)
\arrow[r, "\kappa"] & \Hm_n\left(C \otimes D \right) \arrow[r] &
\displaystyle\bigoplus_{p+q=n-1} \Tor_1^R\left( \Hm_p(C), \Hm_q(D) \right)
\arrow[r] & 0.
\end{tikzcd}
\]
\end{theorem}

% segment-id: o014.li2.chapter3.ext-tor.pf005
\begin{proof}
For every $p\in\Z$, define submodules of $C_p$ by
$B_p:=\Image(d^C_{p+1})$ and $Z_p:=\Ker(d^C_p)$. Regard these as chain
complexes $Z=Z_\bullet$ and $B=B_\bullet$ by setting $d^Z=d^B:=0$.
There is then a short exact sequence of chain complexes
% segment-id: o014.li2.chapter3.ext-tor.g002
\[
\begin{tikzcd}[column sep=large]
0 \arrow[r] & Z \arrow[r] & C \arrow[r, "{d^C=(d^C_p)_p}" inner sep=0.6em]
& B[-1] \arrow[r] & 0.
\end{tikzcd}
\]

The short exact sequence
$0\to Z_p\to C_p\xrightarrow{d^C_p}B_{p-1}\to0$ and the long exact sequence
for $\Tor$ imply that $Z_p$ is also flat. Since $B_{p-1}$ is flat, the
sequence
% segment-id: o014.li2.chapter3.ext-tor.q011
\[
0\to Z_p\otimes D_q\to C_p\otimes D_q\to B_{p-1}\otimes D_q\to0
\]
remains exact. Since direct sums of modules, including infinite direct sums,
preserve exactness, we obtain a short exact sequence of chain complexes
$0\to Z\otimes D\to C\otimes D\xrightarrow{d^C\otimes\identity_D}
B[-1]\otimes D\to0$.

Because $d^Z=d^{B[-1]}=0$, the associated long exact sequence in homology
simplifies, by flatness, to
% segment-id: o014.li2.chapter3.ext-tor.q012
\begin{multline*}
\overbracket{\displaystyle\bigoplus_{p+q=n+1}B_{p-1}\otimes\Hm_q(D)}^{=
\bigoplus_{p+q=n}B_p\otimes\Hm_q(D)}\xrightarrow{\partial_{n+1}}
\displaystyle\bigoplus_{p+q=n}Z_p\otimes\Hm_q(D)\to\Hm_n(C\otimes D)\\
\to\underbracket{\displaystyle\bigoplus_{p+q=n}B_{p-1}\otimes\Hm_q(D)}_{=
\bigoplus_{p+q=n-1}B_p\otimes\Hm_q(D)}\xrightarrow{\partial_n}
\displaystyle\bigoplus_{p+q=n-1}Z_p\otimes\Hm_q(D).
\end{multline*}
It remains only to explain the connecting morphisms $\partial_{n+1}$ and
$\partial_n$. Up to possible signs, they arise from the inclusions
$B_p\hookrightarrow Z_p$. To see this, return to the explicit construction
of the long exact sequence, namely the homological version of
Proposition~\ref{prop:long-exact-sequence-ses}; here one can use a diagram
chase, as in the discussion preceding \cite[Proposition 6.8.6]{Li1}. The
details are left to the reader.

It follows that $\Coker(\partial_{n+1})$ may be identified with
$\bigoplus_{p+q=n}\Hm_p(C)\otimes\Hm_q(D)$, and the induced monomorphism to
$\Hm_n(C\otimes D)$ is $\kappa$.

Finally, $0\to B_p\to Z_p\to\Hm_p(C)\to0$ is a flat resolution of
$\Hm_p(C)$. Hence
$\Ker(\partial_n)\simeq
\bigoplus_{p+q=n-1}\Tor_1^R(\Hm_p(C),\Hm_q(D))$. This gives the asserted
short exact sequence.
\end{proof}

% segment-id: o014.li2.chapter3.ext-tor.p018
Notice that if $R$ is a division ring, then $\kappa$ is necessarily an
isomorphism.

% segment-id: o014.li2.chapter3.ext-tor.p019
Because $D$ plays the role of homology coefficients in topology, the special
case of Theorem~\ref{prop:Kunneth-homology} in which $D$ is concentrated in
degree zero is also called the \emph{universal coefficient theorem}. Results
of this kind have many variants. Moreover, the ring structure supplies
operations on $\Tor$ that are absent for other derived bifunctors. These
structures are better formulated in the language of derived categories or
spectral sequences.
\index{universal coefficient theorem}
